% ============================================================
%  CHAPTER 26: ELECTROMAGNETIC INDUCTION
%  The Physics Codex: A Complete University Foundation
%  Korede Samuel Omotosho (Falcon)
%  Aurikrex Learning Systems, 2026
%
%  File: chapters/part4/ch26_electromagnetic_induction.tex
% ============================================================

\chapter{Electromagnetic Induction}
\label{ch:em_induction}
\minitoc
\newpage

\chapterquote{Nothing is too wonderful to be true, if it
be consistent with the laws of nature; and in such things
as these, experiment is the best test of such
consistency.}{Michael Faraday}

\begin{objectivebox}
By the end of this chapter, you should be able to:
\begin{enumerate}[noitemsep]
  \item Describe the phenomenon of electromagnetic
  induction with experimental evidence
  \item State and apply Faraday's Law of induction
  \item State and apply Lenz's Law
  \item Calculate the EMF induced in a moving
  conductor
  \item Explain and calculate mutual inductance
  and self-inductance
  \item Describe the construction and operation of
  an AC generator
  \item Explain back-EMF in a motor
  \item Describe eddy currents, their causes,
  effects, and uses
  \item Apply electromagnetic induction to practical
  devices including the transformer and generator
\end{enumerate}
\end{objectivebox}

% ============================================================
\section{Intuition: Electricity from Motion}
% ============================================================

In 1831, Michael Faraday made a discovery that would
transform civilisation. He showed that a moving magnet
near a wire could drive a current through it ---
without any battery, without any chemical reaction,
without any contact. Motion itself was generating
electricity.

The device he built was simple: a coil of wire, a
galvanometer, and a bar magnet. When he pushed the
magnet into the coil, the needle deflected. When he
held it still, nothing happened. When he pulled it
out, the needle deflected in the opposite direction.

The key word is \textit{change}. A static magnetic
field produces nothing. A \textit{changing} magnetic
field induces an EMF. This is the entire basis of
modern electrical power generation --- every power
station on Earth, every dynamo, every generator
operates on exactly this principle that Faraday
demonstrated in his laboratory with a magnet and
a coil.

% ============================================================
\section{Faraday's Experiments}
% ============================================================

\begin{diagbox}[Faraday's Induction Experiments]
\begin{center}
\begin{tikzpicture}[scale=0.95, font=\small]

  % === EXPERIMENT 1: Moving magnet into coil ===
  \begin{scope}[shift={(0,5)}]
    \node[font=\bfseries] at (4.0,3.5)
      {Magnet moving into coil $\Rightarrow$ induced current};

    % A two-terminal multi-turn winding; the galvanometer closes its circuit.
    \draw[very thick, gray, decorate,
      decoration={coil, aspect=0.5, segment length=5pt,
      amplitude=5pt}]
      (2.0,1.75) -- (5.0,1.75);
    \draw[thick] (5.0,1.75) -- (6.0,1.75) -- (6.0,1.5);
    \draw[thick] (7.0,1.5) -- (7.0,0.2) -- (2.0,0.2)
      -- (2.0,1.75);

    % Magnet (N end entering from left)
    \fill[codexred!60] (-0.5,1.2) rectangle (0.5,2.3);
    \fill[codexblue!60] (0.5,1.2) rectangle (1.5,2.3);
    \draw[thick] (-0.5,1.2) rectangle (1.5,2.3);
    \node[white, font=\bfseries] at (0.0,1.75) {N};
    \node[white, font=\bfseries] at (1.0,1.75) {S};
    \draw[->, very thick, codexgold] (1.5,1.75) -- (2.0,1.75);
    \node[above, codexgold, font=\small] at (1.75,1.9)
      {$v$};
    \node[font=\tiny, fill=white, inner sep=1pt]
      at (4.2,2.8) {Coil};

    % Galvanometer is inserted between the two circuit terminals.
    \fill[white] (6.5,1.5) circle (0.5);
    \draw[thick] (6.5,1.5) circle (0.5);
    \node[font=\tiny] at (6.5,1.5) {G};
    % Needle deflects right
    \draw[->, thick, codexred] (6.5,1.5) -- (6.9,1.7);

    \node[codexblue, font=\small] at (3.5,-0.8)
      {Needle deflects $\Rightarrow$ induced current};
  \end{scope}

  % === EXPERIMENT 2: Magnet held still ===
  \begin{scope}[shift={(9.5,5)}]
    \node[font=\bfseries] at (3.5,3.5)
      {Magnet stationary $\Rightarrow$ no current};

    % A two-terminal multi-turn winding; the galvanometer closes its circuit.
    \draw[very thick, gray, decorate,
      decoration={coil, aspect=0.5, segment length=5pt,
      amplitude=5pt}]
      (2.0,1.75) -- (5.0,1.75);
    \draw[thick] (5.0,1.75) -- (6.0,1.75) -- (6.0,1.5);
    \draw[thick] (7.0,1.5) -- (7.0,0.2) -- (2.0,0.2)
      -- (2.0,1.75);

    % Magnet static inside coil
    \fill[codexred!60] (1.5,1.2) rectangle (2.5,2.3);
    \fill[codexblue!60] (2.5,1.2) rectangle (3.5,2.3);
    \draw[thick] (1.5,1.2) rectangle (3.5,2.3);
    \node[white, font=\bfseries] at (2.0,1.75) {N};
    \node[white, font=\bfseries] at (3.0,1.75) {S};
    \node[above, gray, font=\small] at (2.5,2.4)
      {stationary};
    \node[font=\tiny, fill=white, inner sep=1pt]
      at (4.2,2.8) {Coil};

    % Galvanometer is inserted between the two circuit terminals.
    \fill[white] (6.5,1.5) circle (0.5);
    \draw[thick] (6.5,1.5) circle (0.5);
    \node[font=\tiny] at (6.5,1.5) {G};
    % Needle vertical (no deflection)
    \draw[-, thick, gray] (6.5,1.5) -- (6.5,1.9);

    \node[gray, font=\small] at (3.5,-0.8)
      {No deflection $\Rightarrow$ no current};
  \end{scope}

  % === EXPERIMENT 3: Magnet withdrawing ===
  \begin{scope}[shift={(0,0)}]
    \node[font=\bfseries] at (4.0,3.5)
      {Magnet withdrawing $\Rightarrow$ opposite current};

    % A two-terminal multi-turn winding; the galvanometer closes its circuit.
    \draw[very thick, gray, decorate,
      decoration={coil, aspect=0.5, segment length=5pt,
      amplitude=5pt}]
      (2.0,1.75) -- (5.0,1.75);
    \draw[thick] (5.0,1.75) -- (6.0,1.75) -- (6.0,1.5);
    \draw[thick] (7.0,1.5) -- (7.0,0.2) -- (2.0,0.2)
      -- (2.0,1.75);

    % Magnet withdrawing
    \fill[codexred!60] (1.5,1.2) rectangle (2.5,2.3);
    \fill[codexblue!60] (2.5,1.2) rectangle (3.5,2.3);
    \draw[thick] (1.5,1.2) rectangle (3.5,2.3);
    \node[white, font=\bfseries] at (2.0,1.75) {N};
    \node[white, font=\bfseries] at (3.0,1.75) {S};
    \draw[->, very thick, codexgold]
      (1.5,1.75) -- (0.8,1.75);
    \node[above, codexgold] at (1.15,1.9) {$v$};
    \node[font=\tiny, fill=white, inner sep=1pt]
      at (4.2,2.8) {Coil};

    % Galvanometer is inserted between the two circuit terminals.
    \fill[white] (6.5,1.5) circle (0.5);
    \draw[thick] (6.5,1.5) circle (0.5);
    \node[font=\tiny] at (6.5,1.5) {G};
    % Needle deflects left (opposite)
    \draw[->, thick, codexblue] (6.5,1.5) -- (6.1,1.7);

    \node[codexred, font=\small] at (3.5,-0.8)
      {Opposite deflection $\Rightarrow$ reversed current};
  \end{scope}

\end{tikzpicture}
\end{center}
\small The \textbf{key observation}: it is the
\textit{change} in magnetic flux that induces EMF ---
not the flux itself.
\end{diagbox}

% ============================================================
\section{Faraday's Law}
% ============================================================

\begin{lawbox}[Faraday's Law of Electromagnetic Induction]
The induced EMF in a circuit is proportional to the
rate of change of magnetic flux linkage through the
circuit:
\[
  \mathcal{E} = -N\frac{d\Phi}{dt}
\]
where:
\begin{itemize}[noitemsep]
  \item $\mathcal{E}$ = induced EMF (volts)
  \item $N$ = number of turns in the coil
  \item $\Phi$ = magnetic flux through one turn
  ($= BA\cos\theta$, in Wb)
  \item $N\Phi$ = total magnetic flux linkage (Wb-turns)
  \item $d\Phi/dt$ = rate of change of flux (Wb\,s$^{-1}$)
\end{itemize}
The magnitude of the induced EMF:
\[
  |\mathcal{E}| = N\frac{\Delta\Phi}{\Delta t}
\]
For a coil with area $A$ in a changing field $B$:
\[
  |\mathcal{E}| = N\frac{\Delta(BA)}{\Delta t}
  = NA\frac{\Delta B}{\Delta t}
  \quad \text{(if area is constant)}
\]
\end{lawbox}

\begin{historybox}[Michael Faraday, 1831]
Michael Faraday (1791--1867) was the son of a
blacksmith and had virtually no formal education ---
he taught himself science by reading books while
working as a bookbinder's apprentice. He became one
of the greatest experimental physicists in history.
His discovery of electromagnetic induction in 1831
led directly to the electric generator and
transformer. He also discovered the laws of
electrolysis, the rotation of polarised light by
magnetic fields, and contributed enormously to
electrochemistry. The unit of capacitance (farad) is
named after him. James Clerk Maxwell, who later
mathematised all of Faraday's experimental findings
into four elegant equations, said Faraday's ideas
were worth more than all the mathematics of the
physicists who came before him.
\end{historybox}

% ============================================================
\section{Lenz's Law}
% ============================================================

\begin{lawbox}[Lenz's Law]
The direction of the induced current is always such
that it \textbf{opposes the change} that caused it.

Equivalently: the induced current creates a magnetic
field that opposes the change in flux through the
circuit.

The negative sign in Faraday's Law ($\mathcal{E}
= -N\,d\Phi/dt$) is a mathematical statement of
Lenz's Law. It reflects conservation of energy ---
if the induced current aided the change, it would
amplify itself indefinitely, creating energy from
nothing, which is impossible.
\end{lawbox}

\begin{diagbox}[Lenz's Law: Direction of Induced Current]
\begin{center}
\begin{tikzpicture}[scale=0.95, font=\small]

  % === MAGNET APPROACHING ===
  \begin{scope}[shift={(0,0)}]
    \node[font=\bfseries] at (3.5,5.5)
      {N pole approaching coil};

    % Coil face
    \draw[very thick] (2.0,1.5) ellipse (1.5cm and 0.5cm);
    \draw[very thick] (2.0,3.5) ellipse (1.5cm and 0.5cm);
    \draw[very thick] (0.5,1.5) -- (0.5,3.5);
    \draw[very thick] (3.5,1.5) -- (3.5,3.5);

    % Magnet approaching from below
    \fill[codexred!60] (1.3,-0.5) rectangle (2.7,0.5);
    \fill[codexblue!60] (1.3,-1.5) rectangle (2.7,-0.5);
    \draw[thick] (1.3,-1.5) rectangle (2.7,0.5);
    \node[white, font=\bfseries] at (2.0,0.0) {N};
    \node[white, font=\bfseries] at (2.0,-1.0) {S};
    \draw[->, very thick, codexgold] (2.0,0.5) -- (2.0,1.3);

    % Increasing flux upward through coil
    \foreach \x in {1.5,2.0,2.5} {
      \draw[->, thick, codexblue!50]
        (\x,0.6) -- (\x,4.0);
    }
    \node[right, codexblue!70, font=\tiny]
      at (3.0,2.5) {Flux $\uparrow$\\(increasing)};

    % Induced current opposes increase:
    % creates field opposing (downward) inside coil
    % By right hand rule: current flows clockwise (viewed from top)
    \draw[very thick, codexred]
      (2.0,3.5) ellipse (1.5 and 0.5);
    \draw[->, very thick, codexred]
      (2.0,3.5) arc(0:270:1.5 and 0.5);
    \node[above, codexred, font=\tiny] at (2.0,4.45)
      {Clockwise current (viewed from top)};
    \node[anchor=west, codexred, font=\tiny, align=left]
      at (4.0,3.2)
      {$B_{ind}$ downward\\(opposes increasing flux)};
  \end{scope}

  % === MAGNET WITHDRAWING ===
  \begin{scope}[shift={(8,0)}]
    \node[font=\bfseries] at (3.5,5.5)
      {N pole withdrawing from coil};

    % Coil
    \draw[very thick] (2.0,1.5) ellipse (1.5cm and 0.5cm);
    \draw[very thick] (2.0,3.5) ellipse (1.5cm and 0.5cm);
    \draw[very thick] (0.5,1.5) -- (0.5,3.5);
    \draw[very thick] (3.5,1.5) -- (3.5,3.5);

    % Magnet withdrawing
    \fill[codexred!60] (1.3,0.5) rectangle (2.7,1.5);
    \fill[codexblue!60] (1.3,-0.5) rectangle (2.7,0.5);
    \draw[thick] (1.3,-0.5) rectangle (2.7,1.5);
    \node[white, font=\bfseries] at (2.0,1.0) {N};
    \node[white, font=\bfseries] at (2.0,0.0) {S};
    \draw[->, very thick, codexgold]
      (2.0,1.3) -- (2.0,0.5);

    % Decreasing flux through coil
    \foreach \x in {1.5,2.0,2.5} {
      \draw[->, thick, codexblue!30, dashed]
        (\x,1.6) -- (\x,4.0);
    }
    \node[right, codexblue!50, font=\tiny]
      at (3.0,2.5) {Flux $\uparrow$\\(decreasing)};

    % Induced current tries to maintain flux:
    % creates upward field, so anticlockwise (from top)
    \draw[very thick, codexblue]
      (2.0,3.5) ellipse (1.5 and 0.5);
    \draw[->, very thick, codexblue]
      (2.0,3.5) arc(0:-270:1.5 and 0.5);
    \node[above, codexblue, font=\tiny] at (2.0,4.45)
      {Anticlockwise current (viewed from top)};
    \node[anchor=west, codexblue, font=\tiny, align=left]
      at (4.0,3.2)
      {$B_{ind}$ upward\\(opposes flux decrease)};
  \end{scope}

\end{tikzpicture}
\end{center}
\end{diagbox}

\begin{examplebox}[26.1]
\stepcounter{workedexample}
A coil of $200$ turns and area $50\ \text{cm}^2$
is placed in a uniform magnetic field of
$0.40\ \text{T}$. The field is then reduced to
zero uniformly in $0.10\ \text{s}$.\\
(a) Calculate the initial magnetic flux through
the coil.\\
(b) Calculate the induced EMF.\\
(c) Using Lenz's Law, state the direction of the
induced current relative to the original flux.
\end{examplebox}

\begin{solutionbox}
\textbf{(a)} Initial flux:
\[
  \Phi_i = BA = 0.40 \times 50\times10^{-4}
  = 2.0\times10^{-3}\ \text{Wb}
\]
Final flux: $\Phi_f = 0$ (field reduced to zero)

\textbf{(b)} Induced EMF:
\[
  |\mathcal{E}| = N\frac{\Delta\Phi}{\Delta t}
  = 200 \times \frac{2.0\times10^{-3} - 0}{0.10}
  = 200 \times 0.020 = 4.0\ \text{V}
\]

\textbf{(c)} The flux is decreasing. By Lenz's Law,
the induced current must \textit{oppose this decrease}
--- it must create a magnetic field in the \textit{same
direction} as the original field, trying to maintain
the flux. Using the right-hand rule, the induced
current flows in the same sense as the current that
would originally have produced the decreasing field.
\end{solutionbox}

% ============================================================
\section{EMF Induced in a Moving Conductor}
% ============================================================

When a conductor of length $L$ moves with velocity $v$
perpendicular to a magnetic field $B$, the free
electrons in the conductor experience a Lorentz force
($F = Bqv$), causing them to accumulate at one end.
This creates a potential difference (EMF) across the
conductor.

\begin{lawbox}[Motional EMF]
For a conductor of length $L$ moving at velocity $v$
perpendicular to a field $B$:
\[
  \mathcal{E} = BLv
\]
This is consistent with Faraday's Law: in time $\Delta t$,
the conductor sweeps out an area $A = Lv\Delta t$,
giving $\Delta\Phi = BLv\Delta t$, so
$\mathcal{E} = \Delta\Phi/\Delta t = BLv$.
\end{lawbox}

\begin{diagbox}[Motional EMF: Conductor Moving in a Field]
\begin{center}
\begin{tikzpicture}[scale=0.82, font=\small]

  % Magnetic field (into page -- crosses)
  \foreach \x in {0.5,1.5,2.5,3.5,4.5,5.5,6.5} {
    \foreach \y in {0.5,1.5,2.5,3.5,4.5} {
      \node[gray!60, font=\small] at (\x,\y) {$\times$};
    }
  }
  \node[gray, font=\small] at (1.5,5.35)
    {$\vect{B}$ (into page)};

  % Conducting rod (vertical, moving right)
  \draw[very thick, codexblue, line width=4pt]
    (3.0,0) -- (3.0,5.0);
  \node[below, codexblue, font=\tiny] at (3.0,-0.35)
    {Rod length $L$};

  % Velocity arrow (right)
  \draw[->, very thick, codexgold, line width=2pt]
    (3.2,2.5) -- (4.5,2.5);
  \node[above, codexgold] at (3.85,2.6) {$v$};

  % Force on positive charges (upward, by F=qv×B)
  \draw[->, thick, codexred]
    (3.0,1.5) -- (3.0,3.5);
  \node[right, codexred, font=\tiny] at (3.9,3.9)
    {$F=qvB$\\on $+$ charges\\(upward)};

  % Charge accumulation
  \node[codexred, font=\bfseries] at (3.0,4.7) {$+$};
  \node[codexred, font=\bfseries] at (3.0,4.4) {$+$};
  \node[codexblue, font=\bfseries] at (3.0,0.6) {$-$};
  \node[codexblue, font=\bfseries] at (3.0,0.3) {$-$};

  % EMF label
  \draw[<->, codexgreen, thick]
    (3.65,0.4) -- (3.65,4.6);
  \node[right, codexgreen] at (3.85,2.5)
    {$\mathcal{E} = BLv$};

  % Rails (horizontal guide)
  \draw[thick, gray] (0,0) -- (7,0);
  \draw[thick, gray] (0,5.0) -- (7,5.0);

  % External circuit
  \draw[thick] (0.5,0) -- (0.5,5.0);
  \fill[white] (0.2,2.2) rectangle (0.8,2.8);
  \draw[thick] (0.2,2.2) rectangle (0.8,2.8);
  \node[font=\tiny] at (0.5,2.5) {R};

  % Current arrow in external circuit
  \draw[->, thick, codexred]
    (0.5,4.8) -- (0.5,3.5);
  \node[left, codexred, font=\tiny] at (0.4,4.2)
    {$I$};

\end{tikzpicture}
\end{center}
\small The rod acts as a source of EMF. Positive
charges are driven upward, so the rod's top is $+$.
Conventional current flows from $+$ to $-$ through
the external circuit (down the left-hand branch), then
through the rod from bottom to top. The rod is the
``battery'' in this circuit.
\end{diagbox}

\begin{examplebox}[26.2]
\stepcounter{workedexample}
An aircraft with a wingspan of $60\ \text{m}$ flies
horizontally at $250\ \text{m\,s}^{-1}$. The vertical
component of the Earth's magnetic field is
$5.0\times10^{-5}\ \text{T}$.\\
(a) Calculate the EMF induced between the wingtips.\\
(b) If the circuit resistance between wingtips (through
the aircraft body) is $0.50\ \Omega$, find the
induced current.\\
(c) Would this current be practically significant?
\end{examplebox}

\begin{solutionbox}
\textbf{(a)} Motional EMF:
\[
  \mathcal{E} = BLv = 5.0\times10^{-5}
  \times 60 \times 250
  = 0.75\ \text{V}
\]

\textbf{(b)} Current:
\[
  I = \frac{\mathcal{E}}{R} = \frac{0.75}{0.50}
  = 1.5\ \text{A}
\]

\textbf{(c)} For an aircraft carrying hundreds of
amperes in its electrical systems, $1.5\ \text{A}$
is negligible. However, the effect is real ---
aircraft do develop charge differences between
wingtips. More significantly, the principle scales
dramatically: conducting bars on ships moving
through the ocean's magnetic field, ocean tidal
water moving through Earth's field, and plasma
flows in the Sun all generate EMFs by exactly
this mechanism.
\end{solutionbox}

% ============================================================
\section{Self-Inductance}
% ============================================================

\begin{definitionbox}[Self-Inductance]
When the current in a coil changes, the changing
magnetic flux through the coil induces an EMF in the
coil itself. This is called \textbf{self-induction}.

The \textbf{self-inductance} $L$ of a coil is defined by:
\[
  \mathcal{E} = -L\frac{dI}{dt}
\]
SI unit: henry (H) where
$1\ \text{H} = 1\ \text{V\,s\,A}^{-1}
= 1\ \text{Wb\,A}^{-1}$.

The self-inductance of a solenoid of $N$ turns,
cross-sectional area $A$, and length $\ell$:
\[
  L = \frac{\mu_0 N^2 A}{\ell} = \mu_0 n^2 A\ell
\]
where $n = N/\ell$ is the turn density.

\textbf{Energy stored in an inductor:}
\[
  W = \frac{1}{2}LI^2
\]
(analogous to $\frac{1}{2}mv^2$ for kinetic energy)
\end{definitionbox}

\begin{keybox}[Self-Inductance and Back-EMF in a Motor]
When a motor spins, the rotating coil generates its
own EMF (by Faraday's Law) that \textbf{opposes}
the applied voltage. This \textbf{back-EMF}
$\mathcal{E}_{back}$ is the same physics as
self-inductance.

When the motor first starts, the back-EMF is zero
(coil not yet moving), so the current is very large:
$I_{start} = V/R_{coil}$.

At full speed, $I = (V - \mathcal{E}_{back})/R_{coil}$.

This is why motors draw very large currents on
startup and why motors burn out when mechanically
stalled (back-EMF $= 0$, only coil resistance limits
current).
\end{keybox}

% ============================================================
\section{Mutual Inductance}
% ============================================================

\begin{definitionbox}[Mutual Inductance]
When two coils are near each other, a changing current
in one coil (primary) induces an EMF in the other
(secondary) through mutual induction.

The \textbf{mutual inductance} $M$ is:
\[
  \mathcal{E}_2 = -M\frac{dI_1}{dt}
\]
SI unit: henry (H).

The transformer is a practical application of mutual
inductance. The mutual inductance between the coils
is maximised by the iron core, which channels all
the flux from the primary through the secondary.
\end{definitionbox}

% ============================================================
\section{The AC Generator}
% ============================================================

\begin{processbox}[How an AC Generator Works]
A rectangular coil of $N$ turns and area $A$ rotates
at angular velocity $\omega$ in a uniform magnetic
field $B$.

The flux through the coil at angle $\theta = \omega t$
from the field:
\[
  \Phi = NBA\cos(\omega t)
\]

By Faraday's Law:
\[
  \mathcal{E} = -N\frac{d\Phi}{dt}
  = NBA\omega\sin(\omega t)
\]
\[
  \boxed{\mathcal{E} = \mathcal{E}_0\sin(\omega t)}
\]
where $\mathcal{E}_0 = NBA\omega$ is the \textbf{peak EMF}.

\textbf{Key features:}
\begin{itemize}[noitemsep]
  \item EMF is zero when the coil is perpendicular
  to the field (flux is maximum but rate of change
  is zero).
  \item EMF is maximum when the coil is parallel to
  the field (flux is zero but rate of change is
  maximum).
  \item \textbf{Slip rings} (not commutator) allow
  the coil to rotate continuously while maintaining
  electrical contact --- this gives AC output.
  \item A \textbf{commutator} converts AC to
  pulsating DC (used in DC generators and motors).
\end{itemize}
\end{processbox}

\begin{diagbox}[AC Generator: Construction and Output]
\begin{center}
\begin{tikzpicture}[scale=0.9, font=\small]

  % === GENERATOR DIAGRAM ===
  \begin{scope}[shift={(0,4)}]
    \node[font=\bfseries] at (3.5,5.3)
      {AC Generator};

    % Magnets
    \fill[codexred!40] (0.0,1.0) rectangle (1.5,4.0);
    \fill[codexblue!40] (5.5,1.0) rectangle (7.0,4.0);
    \draw[thick] (0.0,1.0) rectangle (1.5,4.0);
    \draw[thick] (5.5,1.0) rectangle (7.0,4.0);
    \node[white, font=\bfseries] at (0.75,2.5) {N};
    \node[white, font=\bfseries] at (6.25,2.5) {S};

    % Field lines
    \foreach \y in {1.5,2.5,3.5} {
      \draw[->, thick, codexgold]
        (1.5,\y) -- (5.5,\y);
    }

    % Rotating coil shown as an open loop with two terminals.
    % The loop plane is parallel to B here (maximum induced EMF).
    \draw[very thick, codexblue, line width=2.5pt]
      (4.2,1.0) -- (4.2,4.0) -- (2.8,4.0) -- (2.8,1.0);
    \draw[->, very thick, codexred]
      (4.2,2.0) -- (4.2,3.0);
    \draw[->, very thick, codexblue]
      (2.8,3.0) -- (2.8,2.0);

    % Rotation arrow
    \draw[->, very thick, codexgreen, dashed]
      (3.5,2.5) arc(0:70:1.2);
    \node[codexgreen, font=\small] at (3.0,4.0)
      {$\omega$};

    % Axle
    \fill[gray!50] (3.3,0.3) rectangle (3.7,1.0);
    \fill[gray!50] (3.3,4.0) rectangle (3.7,4.7);

    % Slip rings and separate coil-terminal connections
    \draw[thick] (3.5,0.1) ellipse (0.5 and 0.15);
    \draw[thick] (3.5,-0.3) ellipse (0.5 and 0.15);
    \draw[very thick, codexblue] (2.8,1.0) -- (3.0,0.1);
    \draw[very thick, codexblue] (4.2,1.0) -- (4.0,-0.3);
    \node[below, font=\tiny] at (3.5,-0.55)
      {Slip rings (give AC)};

    % Brushes and external circuit
    \fill[gray!65] (4.0,0.0) rectangle (4.2,0.2);
    \draw[thick] (4.0,0.0) rectangle (4.2,0.2);
    \fill[gray!65] (4.0,-0.4) rectangle (4.2,-0.2);
    \draw[thick] (4.0,-0.4) rectangle (4.2,-0.2);
    \node[above, font=\tiny] at (4.55,0.25) {Brushes};
    \draw[thick] (4.2,0.1) -- (5.0,0.1);
    \draw[thick] (4.2,-0.3) -- (5.0,-0.3);
    \draw[thick] (5.0,0.1) -- (5.0,-1.5);
    \draw[thick] (5.0,-0.3) -- (6.0,-0.3) -- (6.0,-1.5);
    \fill[white] (5.3,-1.5) rectangle (5.7,-2.0);
    \draw[thick] (5.3,-1.5) rectangle (5.7,-2.0);
    \node[right, font=\tiny] at (5.8,-1.75) {External $R$};
    \draw[thick] (5.0,-2.0) -- (5.0,-1.5);
    \draw[thick] (6.0,-2.0) -- (6.0,-1.5);
    \draw[thick] (5.0,-2.0) -- (6.0,-2.0);
  \end{scope}

  % === EMF OUTPUT GRAPH ===
  \begin{axis}[
    at={(220pt,-10pt)},
    width=9cm, height=5cm,
    xlabel={Time ($t$)},
    ylabel={EMF ($\mathcal{E}$)},
    xmin=0, xmax=4*pi,
    ymin=-1.3, ymax=1.5,
    grid=major, grid style={gray!15},
    xtick={0,pi/2,pi,3*pi/2,2*pi,5*pi/2,3*pi,7*pi/2,4*pi},
    xticklabels={0,$T/4$,$T/2$,$3T/4$,$T$,$5T/4$,
      $3T/2$,$7T/4$,$2T$},
    ytick={-1,0,1},
    yticklabels={$-\mathcal{E}_0$,0,$+\mathcal{E}_0$},
    title={AC Output: $\mathcal{E} = \mathcal{E}_0\sin(\omega t)$}
  ]
    \addplot[codexblue, very thick, domain=0:4*pi,
      samples=200]
      {sin(deg(x))};

    % Mark where coil is parallel/perpendicular to field
    \fill[codexred] (axis cs:pi/2,1.0) circle (3pt);
    \node[above, codexred, font=\tiny, align=center,
      text width=1.4cm, fill=white, inner sep=1pt]
      at (axis cs:pi/2,1.05)
      {Coil $\parallel B$\\EMF max};

    \fill[codexblue!60] (axis cs:2*pi,0) circle (3pt);
    \node[above, codexblue, font=\tiny, align=center,
      text width=1.8cm, fill=white, inner sep=1pt]
      at (axis cs:2*pi,1.05)
      {Coil $\perp B$\\$\mathcal{E}=0$};

  \end{axis}

\end{tikzpicture}
\end{center}
\end{diagbox}

\begin{examplebox}[26.3]
\stepcounter{workedexample}
An AC generator has a coil of $80$ turns,
area $120\ \text{cm}^2$, rotating at
$50\ \text{rev\,s}^{-1}$ in a field of
$0.15\ \text{T}$.\\
(a) Calculate the peak EMF.\\
(b) Calculate the RMS EMF.\\
(c) Find the instantaneous EMF at $t = 2.0\ \text{ms}$.
\end{examplebox}

\begin{solutionbox}
$N = 80$, $A = 120\times10^{-4}\ \text{m}^2$,
$f = 50\ \text{Hz}$, $B = 0.15\ \text{T}$

$\omega = 2\pi f = 2\pi \times 50
= 314\ \text{rad\,s}^{-1}$

\textbf{(a)} Peak EMF:
\[
  \mathcal{E}_0 = NBA\omega
  = 80 \times 0.15 \times 120\times10^{-4}
  \times 314
  = 80 \times 0.15 \times 0.012 \times 314
  = 45.2\ \text{V}
\]

\textbf{(b)} RMS EMF:
\[
  \mathcal{E}_{rms} = \frac{\mathcal{E}_0}{\sqrt{2}}
  = \frac{45.2}{\sqrt{2}} = 31.97\ \text{V}
  \approx 32.0\ \text{V}
\]

\textbf{(c)} At $t = 2.0\ \text{ms}
= 2.0\times10^{-3}\ \text{s}$:
\[
  \mathcal{E} = 45.2\sin(314 \times 2.0\times10^{-3})
  = 45.2\sin(0.628)
  = 45.2 \times 0.588
  = 26.6\ \text{V}
\]
\end{solutionbox}

% ============================================================
\section{Eddy Currents}
% ============================================================

\begin{definitionbox}[Eddy Currents]
When a conducting material (not just a wire) moves
through a magnetic field, or when a changing field
passes through it, induced currents circulate within
the bulk of the conductor. These are called
\textbf{eddy currents} (or Foucault currents).

\textbf{Effects of eddy currents:}
\begin{enumerate}[noitemsep]
  \item They dissipate energy as heat in the
  conductor (by $P = I^2R$).
  \item They create a braking force opposing the
  motion (Lenz's Law) --- this is called
  \textbf{magnetic braking}.
\end{enumerate}

\textbf{Reducing unwanted eddy currents:}
The core of a transformer is made of thin,
electrically insulated laminations (sheets) stacked
together. Eddy currents are confined to thin layers,
greatly reducing their magnitude and the associated
energy loss.

\textbf{Useful applications of eddy currents:}
\begin{itemize}[noitemsep]
  \item Induction heating (induction cookers)
  \item Electromagnetic braking (trains, roller
  coasters, braking systems)
  \item Eddy current testing (non-destructive
  testing of metal components for cracks)
  \item Speedometers in vehicles
\end{itemize}
\end{definitionbox}

\begin{diagbox}[Eddy Currents in a Solid vs Laminated Core]
\begin{center}
\begin{tikzpicture}[scale=0.9, font=\small]

  % === SOLID CORE ===
  \begin{scope}[shift={(0,0)}]
    \node[font=\bfseries, codexred] at (2.5,5.0)
      {Solid Core};

    % Solid iron block
    \fill[gray!40] (0.5,0.5) rectangle (4.5,4.5);
    \draw[very thick] (0.5,0.5) rectangle (4.5,4.5);

    % A changing field normal to this cross-section drives eddy currents.
    \foreach \x in {0.9,4.1} {
      \foreach \y in {0.9,4.1} {
        \node[codexblue!70, font=\small] at (\x,\y) {$\otimes$};
      }
    }

    % Large eddy current loops (counterclockwise for increasing inward B)
    \draw[->, very thick, codexred]
      (1.5,1.5) -- (3.5,1.5) -- (3.5,3.5)
      -- (1.5,3.5) -- (1.5,1.5);
    \draw[->, thick, codexred!70]
      (1.0,1.0) -- (4.0,1.0) -- (4.0,4.0)
      -- (1.0,4.0) -- (1.0,1.0);
    \node[codexred, font=\tiny, align=center] at (2.5,-0.85)
      {Large loops; greater $I^2R$ loss};

    \node[below, codexblue, font=\tiny, align=center]
      at (2.5,-0.25) {Increasing $B$ into page};
  \end{scope}

  % === LAMINATED CORE ===
  \begin{scope}[shift={(7,0)}]
    \node[font=\bfseries, codexgreen] at (2.5,5.0)
      {Laminated Core};

    % Laminated iron (thin slices with insulation)
    \foreach \x in {0.5,1.0,1.5,2.0,2.5,3.0,3.5,4.0} {
      \fill[gray!40] (\x,0.5) rectangle (\x+0.3,4.5);
      \draw[thin, white] (\x+0.3,0.5) -- (\x+0.3,4.5);
    }
    \draw[very thick] (0.5,0.5) rectangle (4.5,4.5);

    % The same changing field passes normally through each sheet.
    \foreach \x in {0.7,4.3} {
      \foreach \y in {0.9,4.1} {
        \node[codexblue!70, font=\small] at (\x,\y) {$\otimes$};
      }
    }

    % Small eddy currents (confined to thin sheets)
    \foreach \x in {0.6,1.1,1.6,2.1,2.6,3.1,3.6} {
      \draw[->, thin, codexgreen]
        (\x,1.5) -- (\x+0.2,1.5) -- (\x+0.2,3.5)
        -- (\x,3.5) -- (\x,1.5);
    }
    \node[codexgreen, font=\tiny, align=center] at (2.5,-0.85)
      {Small loops; lower total loss};

    \node[below, codexblue, font=\tiny, align=center]
      at (2.5,-0.25) {Increasing $B$ into page};
  \end{scope}

  % Explanation
  \node[align=center, font=\small] at (5.75,-1.55)
    {Laminations confine eddy currents to thin strips.\\
    Higher resistance path $\Rightarrow$ lower current
    $\Rightarrow$ far less energy loss.};

\end{tikzpicture}
\end{center}
\end{diagbox}

% ============================================================
\section{Chapter Summary}
% ============================================================

\begin{summarybox}
\textbf{Faraday's Law:}
$|\mathcal{E}| = N\,\Delta\Phi/\Delta t$

The induced EMF equals the rate of change of
magnetic flux linkage.

\textbf{Lenz's Law:} Induced current opposes the
change causing it (negative sign in Faraday's Law).
Consequence of energy conservation.

\textbf{Motional EMF:} $\mathcal{E} = BLv$
(conductor of length $L$ moving at velocity $v$
perpendicular to field $B$)

\textbf{Self-inductance:} $\mathcal{E} = -L\,dI/dt$;
$L = \mu_0 N^2 A/\ell$ (solenoid);
Energy: $W = \frac{1}{2}LI^2$.

\textbf{Back-EMF in motor:}
reduces steady-state current; zero at startup
(causes high starting current).

\textbf{AC Generator:}
$\mathcal{E} = NBA\omega\sin(\omega t)
= \mathcal{E}_0\sin(\omega t)$;
peak EMF: $\mathcal{E}_0 = NBA\omega$.
Slip rings give AC; commutator gives DC.

\textbf{Eddy currents:} induced in bulk conductors;
cause braking and heating.
Reduced by lamination in transformer cores.
Used in induction heating, electromagnetic braking.
\end{summarybox}

% ============================================================
\newpage
\begin{exercisebox}[26]

\subsection*{Section A: Multiple Choice}

\textbf{A1.} A coil is placed in a changing magnetic
field. The induced EMF is greatest when:

\mcoption{A}{The flux through the coil is greatest}
\mcoption{B}{The rate of change of flux is greatest}
\mcoption{C}{The flux is zero}
\mcoption{D}{The coil is stationary}
\ansref{Ch26-A1}

\vspace{0.4cm}

\textbf{A2.} A bar magnet is dropped through a
copper tube. Compared with free fall, the magnet
falls:

\mcoption{A}{Faster, due to the magnetic force}
\mcoption{B}{At the same rate, since copper is
non-magnetic}
\mcoption{C}{Slower, because eddy currents create
a braking force}
\mcoption{D}{It stops completely}
\ansref{Ch26-A2}

\vspace{0.4cm}

\textbf{A3.} A straight conductor of length
$0.40\ \text{m}$ moves at $15\ \text{m\,s}^{-1}$
perpendicular to a field of $0.25\ \text{T}$. The
induced EMF is:

\mcoption{A}{$0.10\ \text{V}$}
\mcoption{B}{$0.60\ \text{V}$}
\mcoption{C}{$1.5\ \text{V}$}
\mcoption{D}{$15\ \text{V}$}
\ansref{Ch26-A3}

\vspace{0.4cm}

\textbf{A4.} Lenz's Law is a consequence of:

\mcoption{A}{Coulomb's Law}
\mcoption{B}{Newton's Third Law}
\mcoption{C}{Conservation of energy}
\mcoption{D}{Conservation of charge}
\ansref{Ch26-A4}

\vspace{0.4cm}

\textbf{A5.} The peak EMF of an AC generator is
$\mathcal{E}_0 = NBA\omega$. To double the peak
EMF, without changing $B$ or $A$, you could:

\mcoption{A}{Halve $N$ and double $\omega$}
\mcoption{B}{Double $N$, keeping $\omega$ constant}
\mcoption{C}{Halve $N$ and halve $\omega$}
\mcoption{D}{Double $A$}
\ansref{Ch26-A5}

\vspace{0.4cm}

\textbf{A6.} Transformer cores are laminated in
order to:

\mcoption{A}{Increase the flux linkage}
\mcoption{B}{Reduce hysteresis losses}
\mcoption{C}{Reduce eddy current losses}
\mcoption{D}{Increase the mutual inductance}
\ansref{Ch26-A6}

\vspace{0.4cm}

\textbf{A7.} An inductor of self-inductance $L$
carries current $I$. The energy stored in it is:

\mcoption{A}{$LI$}
\mcoption{B}{$LI^2$}
\mcoption{C}{$\frac{1}{2}LI^2$}
\mcoption{D}{$\frac{1}{2}LI$}
\ansref{Ch26-A7}

\vspace{0.4cm}

\textbf{A8.} The back-EMF in an electric motor
is greatest when:

\mcoption{A}{The motor first starts}
\mcoption{B}{The motor is stalled}
\mcoption{C}{The motor runs at full speed}
\mcoption{D}{The supply voltage is minimum}
\ansref{Ch26-A8}

\vspace{0.4cm}

\textbf{A9.} A coil of 100 turns has flux through
it changing from $0.02\ \text{Wb}$ to
$0.08\ \text{Wb}$ in $0.05\ \text{s}$. The induced
EMF is:

\mcoption{A}{$1.2\ \text{V}$}
\mcoption{B}{$12\ \text{V}$}
\mcoption{C}{$120\ \text{V}$}
\mcoption{D}{$0.12\ \text{V}$}
\ansref{Ch26-A9}

\vspace{0.4cm}

\textbf{A10.} An AC generator uses slip rings
instead of a commutator because:

\mcoption{A}{Slip rings are cheaper}
\mcoption{B}{Slip rings allow the coil to rotate
continuously and deliver alternating current without
reversing it}
\mcoption{C}{Slip rings produce DC output}
\mcoption{D}{Commutators are only used in
transformers}
\ansref{Ch26-A10}

\vspace{0.6cm}

\subsection*{Section B: Theory and Calculations}

\textbf{B1.} A flat circular coil of 150 turns
and radius $4.0\ \text{cm}$ is placed with its
plane perpendicular to a uniform magnetic field.
The field changes uniformly from $0.60\ \text{T}$
to $0.10\ \text{T}$ in $0.20\ \text{s}$.\\[4pt]
(a) Calculate the initial and final magnetic flux
through the coil.\\
(b) Calculate the magnitude of the induced EMF.\\
(c) If the coil has resistance $5.0\ \Omega$, find
the induced current.\\
(d) State the direction of the induced current
relative to the decreasing field (use Lenz's Law).\\
(e) If instead of the field changing, the coil is
rotated $90°$ in $0.20\ \text{s}$ while the field
remains at $0.60\ \text{T}$, calculate the average
EMF induced.
\hfill\ansref{Ch26-B1}

\vspace{0.4cm}

\textbf{B2.} A pair of horizontal conducting rails
are $50\ \text{cm}$ apart in a vertical magnetic
field of $0.80\ \text{T}$ (perpendicular to the
rails). A conducting rod of mass $30\ \text{g}$ and
resistance $0.40\ \Omega$ lies across the rails and
is free to slide. The rails have negligible resistance.\\[4pt]
(a) The rod is pushed at a constant speed of
$2.5\ \text{m\,s}^{-1}$ along the rails.
Calculate the induced EMF and current.\\
(b) Calculate the force required to push the rod
at this constant speed (the magnetic braking force
equals the applied force at constant velocity).\\
(c) Calculate the power input by the pushing force
and the power dissipated in the rod. Comment on
the relationship between them.\\
(d) If the external force is removed, the rod
decelerates due to the braking force. Write the
equation of motion ($F = ma$) and explain why the
deceleration decreases as the rod slows down.\\
(e) This system is a linear generator. At what
rod speed would it generate $1.0\ \text{V}$ across
the rails?
\hfill\ansref{Ch26-B2}

\vspace{0.4cm}

\textbf{B3.} An AC generator consists of a coil
with 500 turns of area $80\ \text{cm}^2$ rotating
in a field of $0.20\ \text{T}$. The coil makes
$25$ complete revolutions per second.\\[4pt]
(a) Calculate the angular frequency $\omega$.\\
(b) Calculate the peak EMF.\\
(c) Calculate the RMS EMF.\\
(d) The coil has resistance $2.0\ \Omega$ and is
connected to an external load of $50\ \Omega$.
Calculate the peak current, RMS current, and
average power delivered to the load.\\
(e) Explain why the coil requires a torque to keep
it rotating at constant speed against the
electromagnetic braking force.\\
(f) The rotational speed is increased to
$50\ \text{rev\,s}^{-1}$. Find the new peak EMF.
\hfill\ansref{Ch26-B3}

\vspace{0.4cm}

\textbf{B4.} A solenoid has $600$ turns, length
$25\ \text{cm}$ and cross-sectional area
$8.0\ \text{cm}^2$.
($\mu_0 = 4\pi\times10^{-7}\ \text{T\,m\,A}^{-1}$)\\[4pt]
(a) Calculate the self-inductance of the solenoid.\\
(b) The current in the solenoid increases uniformly
from $0$ to $4.0\ \text{A}$ in $50\ \text{ms}$.
Calculate the back-EMF.\\
(c) Calculate the energy stored in the solenoid
when the current is $4.0\ \text{A}$.\\
(d) A second identical solenoid is wound coaxially
around the first and the mutual inductance between
them is $M = 0.72\ \mu\text{H}$. When the current
in solenoid 1 changes at $200\ \text{A\,s}^{-1}$,
calculate the induced EMF in solenoid 2.\\
(e) Explain how an iron core increases the
self-inductance of the solenoid, referring to the
formula $L = \mu_r\mu_0 N^2A/\ell$.
\hfill\ansref{Ch26-B4}

\vspace{0.4cm}

\textbf{B5.} An electric motor runs on a $240\ \text{V}$
DC supply. Its armature coil has resistance
$2.0\ \Omega$. When running at full speed, it draws
a current of $3.0\ \text{A}$.\\[4pt]
(a) Calculate the back-EMF of the motor at full
speed.\\
(b) Calculate the power input, power dissipated as
heat in the coil, and mechanical power output.\\
(c) Calculate the efficiency of the motor.\\
(d) When the motor is first switched on (before it
starts spinning), calculate the current drawn.
Explain why this is potentially damaging.\\
(e) A starter resistor of $18\ \Omega$ is connected
in series with the motor at startup. Calculate the
current on startup with the starter resistor and
the power dissipated in it. At what fraction of
full speed would you remove the starter resistor,
given that you want the motor current not to exceed
$6.0\ \text{A}$?
\hfill\ansref{Ch26-B5}

\end{exercisebox}
